\textbf{Definition (recursive rule).} Define $f(1)=f(2)=1$. For $n\ge 3$, attempt to define
\[
 f(n) = f\big(n-f(n-1)\big)+f\big(n-f(n-2)\big),
\]
provided that the arguments $n-f(n-1)$ and $n-f(n-2)$ are positive integers for which $f(\cdot)$ has already been defined.

\textbf{Well-definedness.} Say the sequence is \emph{well-defined up to $N$} if for every $3\le n\le N$, both indices $n-f(n-1)$ and $n-f(n-2)$ lie in $\{1,2,\dots,n-1\}$.

\textbf{Questions.}
\begin{itemize}
\item[(Q1)] Is $f(n)$ well-defined for all $n\in\mathbb{N}$?
\item[(Q2)] Does the range $\{f(n): n\in\mathbb{N}\}$ omit infinitely many positive integers?
\item[(Q3)] Describe the growth/behaviour of $f(n)$.
\end{itemize}
